\subsection{Alcances}

Los alcances del proyecto comprenden:

\begin{itemize}
  \item El proyecto se enfocará en el desarrollo de un prototipo jugable de videojuego de simulación gamificada, orientado al apoyo del aprendizaje práctico de la ciberseguridad en estudiantes de Ingeniería Informática e Ingeniería de Sistemas.
  \item Se abordará la representación de escenarios controlados de amenazas digitales fundamentales, tales como phishing, malware e ingeniería social, mediante situaciones interactivas que exijan al jugador identificar riesgos y seleccionar medidas de respuesta.
  \item Se implementarán mecánicas de juego basadas en toma de decisiones, donde el jugador podrá aplicar contramedidas defensivas como capacitación, actualización de sistemas, uso de firewalls, políticas de seguridad y otras acciones preventivas o correctivas.
  \item Se desarrollará una progresión de escenarios en la que la dificultad y la complejidad conceptual aumenten gradualmente, permitiendo que el estudiante avance desde situaciones básicas hacia problemas con mayor nivel de análisis.
  \item El videojuego incluirá retroalimentación inmediata sobre las decisiones realizadas por el jugador, mostrando consecuencias como amenazas contenidas, daño mitigado, recursos protegidos, pérdidas generadas o puntuación de efectividad.
  \item Se incorporará un módulo de seguimiento de progreso que registre información básica del desempeño del jugador, como escenarios completados, decisiones críticas, puntuación obtenida, tiempo de reacción y patrones generales de defensa.
  \item Se realizarán pruebas de usabilidad y funcionamiento con usuarios, con el propósito de verificar que el prototipo sea comprensible, funcional y adecuado como herramienta complementaria de aprendizaje.
\end{itemize}
